% PAGES: 117-118

Consider the following seven securities:\footnote{More choices of
securities are discussed in exercises at the end of the chapter.}
\begin{equation}
\begin{aligned}
&\text{P1200;} \\
&\text{P1300;} \\
&\text{P1400;} \\
&\text{C1400;} \\
&\text{C1450;} \\
&\text{C1550;} \\
&\text{C1600.}
\end{aligned}
\end{equation}

Note that the options from (3.43) are out--of--the--money puts and
out--of--the--money calls,\footnote{The choice of out--of--the--money
options is driven by the fact that options market makers only trade
out--of--the--money options, which have the same implied volatilities as
corresponding in--the--money options. For example, the trading volume of
the out--of--the--money options (C1400 through C1600 and P1175 through
P1350) from Table 3.1 is $32,359$, while the trading volume of the
in--the--money options (C1175 through C1350 and P1400 through P1600 is
$3,134$, about 10 times smaller.} plus a put option (P1400) and a call
option (C1400) struck around at--the--money; recall that the spot price
of the index was $1,370$. For details on at--the--money,
out--of--the--money, and in--the--money options, see Section 10.3.

The price vector of the securities from (3.43) at time $t_0$ is the
vector of their mid prices from Table 3.1, i.e.,
\begin{equation}
S_{t_0} \ = \ \begin{pmatrix} 51.55 \\ 77.70 \\ 116.55 \\ 71.10 \\ 47.25 \\ 15.80 \\ 7.90 \end{pmatrix}.
\end{equation}

Seven possible states of the index price on 12/22/2012 are obtained as
follows: There are six option strikes corresponding to the securities
above, i.e.,
\[
1200, 1300, 1400, 1450, 1550, 1600.
\]

The five midpoints
\[
1250, 1350, 1425, 1500, 1575
\]
between the strikes above will give the following five states:
\begin{alignat*}{4}
\omega^2 &: \{S(\tau) = 1250\}; &\qquad \omega^5 &: \{S(\tau) = 1500\}; \\
\omega^3 &: \{S(\tau) = 1350\}; &\qquad \omega^6 &: \{S(\tau) = 1575\}; \\
\omega^4 &: \{S(\tau) = 1425\}, & &
\end{alignat*}
where $S(\tau)$ denotes the value of the S\&P 500 index on 12/22/2012.
Below the lowest strike and above the highest strike we choose the
following two states:\footnote{The choice of the two states below the
smallest option strike and above the largest strike could influence
whether the one period market model is arbitrage--free, and the overall
performance of the risk--neutral pricing of the other derivative
securities.}
\[
\omega^1 : \{S(\tau) = 1000\}; \qquad \omega^7 : \{S(\tau) = 1700\}.
\]

Recall from (10.90) and (10.91) that the payoffs at maturity of call and
put options are, respectively,
\begin{align}
C(T) &= \max(S(T)-K,0) = \begin{cases} S(T)-K, & \text{if } S(T) > K; \\ 0, & \text{if } S(T) \le K. \end{cases} \\
P(T) &= \max(K-S(T),0) = \begin{cases} 0, & \text{if } S(T) \ge K; \\ K-S(T), & \text{if } S(T) < K. \end{cases}
\end{align}

Then, the payoff matrix of this one period market model is
\begin{equation}
M_\tau \ = \
\begin{pmatrix}
200 & 0   & 0  & 0  & 0   & 0   & 0   \\
300 & 50  & 0  & 0  & 0   & 0   & 0   \\
400 & 150 & 50 & 0  & 0   & 0   & 0   \\
0   & 0   & 0  & 25 & 100 & 175 & 300 \\
0   & 0   & 0  & 0  & 50  & 125 & 250 \\
0   & 0   & 0  & 0  & 0   & 25  & 150 \\
0   & 0   & 0  & 0  & 0   & 0   & 100
\end{pmatrix}.
\end{equation}

For example, $M_\tau(3,3)$ is the payoff of the third security P1400,
i.e., of the put option with strike 1400, if state $\omega^3$ occurs,
i.e., if $S(\tau) = 1350$, and from (3.46) we find that
\[
M_\tau(3,3) \ = \ \max(1400-S(\tau),0) \ = \ \max(1400-1350,0) \ = \ 50.
\]

Similarly, $M_\tau(5,6)$ is the payoff of the fifth security C1450,
i.e., of the call option with strike 1450, if state $\omega^6$ occurs,
i.e., if $S(\tau) = 1575$, and from (3.45) we find that
\[
M_\tau(5,6) \ = \ \max(S(\tau)-1450,0) \ = \ \max(1575-1450,0) \ = \ 125.
\]

We now show that this model is complete and arbitrage--free.

Note that the matrix $M_\tau$ is made of a $3 \times 3$ lower triangular
block and a $4 \times 4$ upper triangular block. Thus, its determinant
$\det(M_\tau)$ is the product of the diagonal entries of $M_\tau$, and
therefore nonzero; see Section 10.1.1. It follows that the matrix
$M_\tau$ is nonsingular, see Theorem 1.2, and, from Lemma 3.4, we
conclude that this market model is complete.

To decide whether this market model is arbitrage--free, we must solve
the linear system $M_\tau Q = S_{t_0}$, where $M_\tau$ and $S_{t_0}$ are
given by (3.47) and (3.44), respectively, and check whether all the
entries of the vector $Q$ are positive.

Note that, due to the block structure of the payoff matrix $M_\tau$,
solving the linear system $M_\tau Q = S_{t_0}$ amounts to a forward
substitution corresponding to the $3 \times 3$ lower triangular block of
the matrix $M_\tau$ for computing the first 3 entries of $Q$, and a
backward substitution corresponding to the $4 \times 4$ upper triangular
block of $M_\tau$ for computing the last 4 entries of $Q$, i.e.,
\[
Q \ = \ \left(
\begin{array}{l}
\text{forward\_subst}\!\left(
\begin{pmatrix} 200 & 0 & 0 \\ 300 & 50 & 0 \\ 400 & 150 & 50 \end{pmatrix},
\begin{pmatrix} 51.55 \\ 77.70 \\ 116.55 \end{pmatrix}
\right) \\[3em]
\text{backward\_subst}\!\left(
\begin{pmatrix} 25 & 100 & 175 & 300 \\ 0 & 50 & 125 & 250 \\ 0 & 0 & 25 & 150 \\ 0 & 0 & 0 & 100 \end{pmatrix},
\begin{pmatrix} 71.10 \\ 47.25 \\ 15.80 \\ 7.90 \end{pmatrix}
\right)
\end{array}
\right).
\]
